\documentclass[12pt,a4paper]{article}

% ── Packages ──────────────────────────────────────────────────────────────────
\usepackage[utf8]{inputenc}
\usepackage[T1]{fontenc}
\usepackage{amsmath,amssymb,amsthm}
\usepackage{mathrsfs}
\usepackage{geometry}
\usepackage{booktabs}
\usepackage{array}
\usepackage{longtable}
\usepackage{hyperref}
\usepackage{xcolor}
\usepackage{graphicx}
\usepackage{setspace}
\usepackage{enumitem}
\usepackage{caption}

\geometry{margin=1in}
\onehalfspacing

% ── Theorem environments ───────────────────────────────────────────────────────
\newtheorem{theorem}{Theorem}[section]
\newtheorem{proposition}[theorem]{Proposition}
\newtheorem{lemma}[theorem]{Lemma}
\newtheorem{corollary}[theorem]{Corollary}
\newtheorem{remark}{Remark}[section]
\newtheorem{definition}{Definition}[section]

% ── Custom commands ────────────────────────────────────────────────────────────
\newcommand{\Cl}{\mathrm{Cl}}
\newcommand{\SU}{\mathrm{SU}}
\renewcommand{\Re}{\operatorname{Re}}
\newcommand{\Tr}{\operatorname{Tr}}
\newcommand{\diag}{\operatorname{diag}}
\newcommand{\embed}{\operatorname{embed}}
\newcommand{\HH}{\mathbb{H}}
\newcommand{\RR}{\mathbb{R}}
\newcommand{\CC}{\mathbb{C}}
\newcommand{\OO}{\mathbb{O}}
\newcommand{\Id}[1]{I_{#1}}
\newcommand{\Pew}[1]{P^{#1}_{\mathrm{EW}}}
\newcommand{\GW}[1]{G_W^{(#1)}}
\newcommand{\fEW}[1]{f_{#1}^{\mathrm{EW}}}
\newcommand{\sigx}{\sigma_x}
\newcommand{\sigy}{\sigma_y}
\newcommand{\sigz}{\sigma_z}
\newcommand{\brac}[2]{\left[#1,\,#2\right]}

% ── Hyperref setup ─────────────────────────────────────────────────────────────
\hypersetup{
	colorlinks=true,
	linkcolor=blue!70!black,
	citecolor=green!50!black,
	urlcolor=blue!60!black
}
% Suppress math-in-PDF-bookmark warnings from hyperref
\pdfstringdefDisableCommands{%
	\def\Cl#1{Cl(#1)}%
	\def\SU#1{SU(#1)}%
	\def\HH{H}%
	\def\RR{R}%
	\def\CC{C}%
	\def\OO{O}%
	\def\mathbb#1{#1}%
	\def\mathrm#1{#1}%
	\def\textsubscript#1{_#1}%
}

% ══════════════════════════════════════════════════════════════════════════════
\begin{document}
	
	\title{\textbf{Per-Configuration Equality of the T\textsubscript{3} Grade\\ 
			Channels of the SU(2) Link Correlator\\
			in the $\Cl(6)/\HH$ Framework}\\[6pt]}
	
	\author{Todd M. Firestone}
	
	\date{Keywords: Clifford algebras, SU(2) gauge theory, lattice field theory,
		grade decomposition, pseudoreality, T\textsubscript{3} link correlator}
	
	\maketitle

	
	% ── Abstract ──────────────────────────────────────────────────────────────────
	\begin{abstract}
		We decompose the SU(2) pure-gauge link correlator into two $T_3$ grade
		channels within the $\Cl(2)\cong\HH$ algebraic framework of this series.
		The channel projectors $\Pew{\pm}$, defined as the Lagrange spectral
		projectors of the SU(2) Cartan generator
		$T_3 = \tfrac{1}{2}(I_2\otimes\sigma_z\otimes I_2)$ acting on the
		eight-dimensional representation space, split the correlator into two pieces
		of equal dimension ($4+4$).  Our central result is an exact, per-configuration
		identity: for every SU(2) gauge configuration the two channel correlators are
		equal, $G_W^{(+)}(t)=G_W^{(-)}(t)$.  This is an algebraic consequence of SU(2)
		pseudoreality --- the diagonal elements of any $U\in\SU(2)$ satisfy
		$U_{11}=\overline{U_{00}}$ --- and holds configuration by configuration, not
		merely on ensemble average.  We prove the identity and verify it to machine
		precision on thermalized SU(2) ensembles at $\beta\in\{4.5,6.0\}$.  We delimit
		the result carefully.  The channel grade fractions
		$\fEW{\pm}(t)=G_W^{(\pm)}(t)/G_W(t)$ equal $1/2$ only in the free field
		(all links $U_\mu=I_2$); in the interacting theory the cross term
		$\Delta=G_W-G_W^{(+)}-G_W^{(-)}$ is nonzero, so the fractions are not pinned
		to $1/2$.  Moreover, the channel correlators and their cross term are
		gauge-variant: by Elitzur's theorem their ensemble averages vanish in the
		absence of gauge fixing and are not physical observables.  The surviving
		content is therefore the per-configuration channel equality --- a structural
		feature of the framework rooted in pseudoreality, with no analogue in the
		SU(3) color sector.  This paper does not address the Higgs sector (excluded by
		the Derivability Limit, stated in \S2.4), chiral fermions (excluded by
		Nielsen--Ninomiya), or the L/R handedness of the SU(2) gauge group (the L/R
		selection requires external physical input not provided by the algebra).
	\end{abstract}
	
	
	% ══════════════════════════════════════════════════════════════════════════════
	\section{Introduction}
	
\subsection{Context and Question}

This series develops an algebraic framework for embedding the Standard Model gauge
groups within a tensor product of division algebras
$\RR\otimes\CC\otimes\HH\otimes\OO$~\cite{P1,P2}.  Paper~2~\cite{P2} studied the
grade-decomposed QCD pion correlator: the color-grade projectors of the $\Cl(6)$
representation space split the correlator into channels of dimension
$3,3,1,1$, and in the free Wilson--Dirac theory the corresponding fractions take
the kinematic values $\dim(r)/8$.

The present paper asks the analogous structural question in the electroweak sector.
Embedding SU(2) in the eight-dimensional representation space, the Cartan generator
$T_3$ has the two-valued spectrum $\{\pm\tfrac12\}$ and its spectral projectors
$\Pew{\pm}$ split the SU(2) link correlator into two channels of equal dimension
($4+4$).  We ask: \emph{what, if anything, is exact about this $T_3$ grade
	decomposition?}

The answer is a single, sharp identity.  For every SU(2) gauge configuration the
two channel correlators are equal,
\[
G_W^{(+)}(t) = G_W^{(-)}(t),
\]
a per-configuration consequence of SU(2) pseudoreality.  This is the central result
of the paper.  We also show, and emphasize as a corrective to a natural but false
expectation, that the channel \emph{fractions} $\fEW{\pm}(t)$ equal $1/2$ only in
the free field; in the interacting theory a nonzero cross term spoils the equality
of the fractions even though the channel correlators themselves remain equal.

The computation lies entirely in the bosonic sector, sidestepping the
Nielsen--Ninomiya obstruction~\cite{NN1,NN2} that prevents the inclusion of chiral
fermions.
	
	\begin{remark}[L/R ambiguity]
		The generators $T_1,T_2,T_3$ embed SU(2) in the second tensor factor of the
		8-dimensional space via $\Id{2}\otimes\sigma_{x,y,z}/2\otimes\Id{2}$.  This
		embedding does not, by itself, select $\SU(2)_L$ over $\SU(2)_R$; distinguishing
		the two requires external physical input (such as coupling to left-handed matter
		fields) not provided by the algebra.  Throughout this paper, ``SU(2)'' means the
		abstract gauge group, with no claim about chirality.
	\end{remark}
	
\subsection{Prior Work and Context}

SU(2) pure gauge lattice theory has a long history~\cite{K1979,BC1980,L1998}.
In the action normalization used here (\S\ref{sec:lattice}), the SU(2)
deconfinement transition lies near $\beta\approx 4.6$; the two couplings studied,
$\beta\in\{4.5,6.0\}$, straddle this value, so the ensembles are not uniformly in
the weak-coupling phase.  The novelty of the present work lies not in the SU(2)
gauge dynamics themselves, but in the $T_3$ grade decomposition of the link
correlator and the exact channel identity it exposes.

The $\Cl(2)\cong\HH$ isomorphism identifies the quaternion algebra $\HH$ with the
two-dimensional Clifford algebra, which provides the natural algebraic language for
SU(2) gauge theory.  The channel projectors $\Pew{\pm}$ arise as spectral projectors
of $T_3$, a generator with eigenspectrum $\{\pm1/2\}$ only, reflecting the symmetric
doublet structure of SU(2).  This is distinct from the SU(3) color sector of
Paper~2, whose projectors have unequal dimensions ($3+3+1+1$): the equal-dimension
$T_3$ split, combined with the pseudoreality of SU(2), is what yields the exact
per-configuration channel equality $G_W^{(+)}=G_W^{(-)}$ that constitutes the main
result.
	
\subsection{Scope}

The following are \textbf{explicitly out of scope}:

\begin{enumerate}[leftmargin=*,label=\arabic*.]
	\item \textbf{No Higgs sector.}  The Derivability Limit (stated in \S2.4)
	establishes that the division-algebra structure cannot constrain the Higgs
	potential $V(\phi)$.  The Higgs sector is the subject of a separate programme.
	\item \textbf{No chiral fermions.}  The Nielsen--Ninomiya theorem~\cite{NN1,NN2}
	prohibits exact chiral symmetry on the lattice with a local, hermitian, gapless
	Dirac operator.
	\item \textbf{No L/R distinction.}  The embedding does not select a chirality.
	The results apply to the abstract SU(2) gauge group, with the explicit caveat
	that L/R assignment requires external input.
	\item \textbf{No ensemble-averaged channel observables.}  The channel
	correlators $G_W^{(\pm)}$ and their cross term are gauge-variant; their ensemble
	averages vanish by Elitzur's theorem in the absence of gauge fixing (\S\ref{sec:gauge}).
	The paper's result is the \emph{per-configuration} channel identity, not a
	statement about expectation values.
	\item \textbf{No dynamical symmetry-breaking signal.}  No claim about
	confinement or electroweak symmetry breaking is made.
\end{enumerate}
	
	% ══════════════════════════════════════════════════════════════════════════════
	\section{Algebraic Framework}
	
	\subsection{SU(2) Generators in the 8-Dimensional Space}
	
	We embed SU(2) in the 8-dimensional $\Cl(6)$ representation space through the
	generators
	\begin{equation}
		T_1 = \tfrac{1}{2}(\Id{2}\otimes\sigx\otimes\Id{2}),\quad
		T_2 = \tfrac{1}{2}(\Id{2}\otimes\sigy\otimes\Id{2}),\quad
		T_3 = \tfrac{1}{2}(\Id{2}\otimes\sigz\otimes\Id{2}),
	\end{equation}
	where $\sigma_{x,y,z}$ are the Pauli matrices and all generators act on the
	\emph{second} tensor factor of the decomposition $\CC^2\otimes\CC^2\otimes\CC^2$.
	These satisfy the SU(2) Lie algebra $\brac{T_i}{T_j}=i\varepsilon_{ijk}T_k$, and
	the quadratic Casimir $T^2=T_1^2+T_2^2+T_3^2=\tfrac{3}{4}\Id{8}$, consistent
	with the spin-$\tfrac{1}{2}$ representation on each of the two $\CC^2$ factors.
	Both identities are verified at machine precision (checks A1 and A2 in
	Table~\ref{tab:algebra}).
	
	The generator $T_3$ has eigenvalues
	$\{+\tfrac{1}{2},+\tfrac{1}{2},+\tfrac{1}{2},+\tfrac{1}{2},
	-\tfrac{1}{2},-\tfrac{1}{2},-\tfrac{1}{2},-\tfrac{1}{2}\}$,
	each with multiplicity~4 (check~A3).  This equal-multiplicity structure is a
	direct consequence of the SU(2) algebra: the outer
	$\CC^2\otimes\CC^2$ factor carries no SU(2) quantum number, contributing a factor
	of~4 to each eigenspace.
	
	The \textbf{link embedding} that places a $2\times2$ SU(2) link matrix into the
	8-dimensional space is
	\begin{equation}
		U_\mu^{(8)}(x) = \Id{2}\otimes U_\mu(x)\otimes\Id{2},
	\end{equation}
	consistent with the second-tensor-factor placement of the generators.  Embedding
	unitarity,
	$\|\,(\Id{2}\otimes U\otimes\Id{2})(\Id{2}\otimes U\otimes\Id{2})^\dagger
	-\Id{8}\,\|_F < 1.3\times10^{-15}$,
	is confirmed for 20 random SU(2) matrices (check~B7).
	
	\subsection{Electroweak Projectors}
	
	The $T_3$ eigenspaces define the electroweak projectors via Lagrange interpolation:
	\begin{equation}
		\Pew{+} = \frac{\Id{8}+2T_3}{2}\quad(T_3=+{\tfrac{1}{2}}\ \text{eigenspace, dim}=4),
	\end{equation}
	\begin{equation}
		\Pew{-} = \frac{\Id{8}-2T_3}{2}\quad(T_3=-{\tfrac{1}{2}}\ \text{eigenspace, dim}=4).
	\end{equation}
	Eight checks pass at machine precision (checks B1--B6, B8,
	Table~\ref{tab:algebra}): completeness ($\Pew{+}+\Pew{-}=\Id{8}$), orthogonality
	($\Pew{+}\Pew{-}=0$), idempotence ($(\Pew{\pm})^2=\Pew{\pm}$), hermiticity, and
	equal trace ($\Tr\Pew{\pm}=4$).  The key structural fact ---
	$\Tr\Pew{+}=\Tr\Pew{-}=4$ --- reflects the equal-dimension decomposition $8=4+4$.
	This differs from the color sector of Paper~2, where the four projectors have
	dimensions $3,3,1,1$ (an asymmetric decomposition, $8=3+3+1+1$).
	
	Check~B8 verifies that the EW projectors commute with the color Cartan projectors
	of Paper~2: $[P_{\mathrm{color}}^{(r)},\Pew{s}]=0$ to machine precision,
	confirming that the color and EW gradings are compatible decompositions of the
	same 8-dimensional space.
	
	\subsection{Constraint from the \texorpdfstring{$T_3$}{T3} Eigenspectrum}
	
	The $T_3$ eigenspectrum $\{\pm1/2\}$ (with no zero eigenvalue) has a direct
	physical implication: the 8-dimensional representation space contains no singlet
	under $\SU(2)_L$.  Every state is assigned to the doublet representation.
	Right-handed fermions (which carry $T_3=0$ in the Standard Model) cannot be
	accommodated within this representation.  This confirms the bosonic and
	doublet-sector restriction of the present calculation.  The projector algebra used
	below is established by the machine-precision checks of Table~\ref{tab:algebra}.
	

	\subsection{The Derivability Limit}
	
	For completeness we record the proposition, established earlier in this series,
	that bounds what the algebra can determine.
	
	\begin{proposition}[Derivability Limit]
		The division-algebra structure $\RR\otimes\CC\otimes\HH\otimes\OO$ and its
		$\Cl(6)$ representation fix the gauge group embeddings, the grade decompositions,
		and the representation content, but do not fix the coefficients of any
		gauge-invariant potential built from matter fields. In particular, the Higgs
		potential $V(\phi)$ --- its mass term and self-coupling --- is not constrained by
		the algebra and must be supplied as external physical input.
	\end{proposition}
	
	\noindent
	The Derivability Limit is the reason the Higgs sector is excluded from this paper
	(\S1.3): the present results concern only the algebraically fixed grade structure of
	the pure-gauge link correlator, not any quantity requiring a scalar potential.
	
	% ══════════════════════════════════════════════════════════════════════════════
	\section{Lattice Implementation}
	\label{sec:lattice}
	
	\subsection{SU(2) Wilson Gauge Action}
	
	We simulate 4-dimensional SU(2) pure Yang--Mills theory on a hypercubic lattice
	of linear extent $N=6$ and volume $V=N^4=1296$ with periodic boundary conditions.
	The Wilson gauge action is
	\begin{equation}
		S_{\SU(2)} = \frac{\beta}{2}\sum_{x,\,\mu<\nu}
		\left[1-\frac{1}{2}\Re\Tr P_{\mu\nu}(x)\right],
	\end{equation}
	where the plaquette
	$P_{\mu\nu}(x)=U_\mu(x)\,U_\nu(x+\hat\mu)\,U_\mu^\dagger(x+\hat\nu)\,U_\nu^\dagger(x)$
	and $U_\mu(x)\in\SU(2)$.  With the action normalization above
	($S=\tfrac{\beta}{2}\sum_P[1-\tfrac12\Re\Tr P]$), the coupling $\beta$ is twice the
	standard SU(2) Wilson coupling $\beta_{\mathrm{std}}=4/g^2$, so the deconfinement
	transition at $\beta_{\mathrm{std}}\approx2.30$ corresponds to $\beta\approx4.6$ in
	our convention.  We study $\beta\in\{4.5,6.0\}$ ($\beta_{\mathrm{std}}\in\{2.25,3.0\}$),
	which straddle the transition: $\beta=4.5$ lies just below it (confined side) and
	$\beta=6.0$ just above (deconfined side).  None of our results depend on the phase,
	as will be made explicit.
	
	The plaquette observable is measured as
	\begin{equation}
		\langle P\rangle = \frac{1}{V}\cdot\frac{1}{6}\sum_{x,\mu<\nu}
		\frac{\Re\Tr[P_{\mu\nu}(x)]}{2}
	\end{equation}
	and reported in Table~\ref{tab:plaquette}.
	
	\subsection{Even-Odd Metropolis Update}
	
	Link updates employ the even-odd checkerboard Metropolis algorithm established in
	Paper~2~\S3.2.  Sites are partitioned by the parity of the coordinate sum,
	$\mathfrak{p}(x)=(\sum_\mu x_\mu)\bmod 2\in\{0,1\}$.  Within each parity class
	and each direction $\mu$, all links are proposed and accepted or rejected in
	parallel.
	
	For each link, the Metropolis proposal is generated as
	\begin{equation}
		U_\mu^{\mathrm{new}}(x) = R\cdot U_\mu^{\mathrm{old}}(x),\qquad
		R = \exp\!\bigl[i\varepsilon\,(r_1\sigx+r_2\sigy+r_3\sigz)\bigr],
	\end{equation}
	where $r_1,r_2,r_3\sim\mathcal{N}(0,1)$ are standard normal variates and
	$\varepsilon$ is the step-size parameter.  The proposal is projected back to SU(2)
	by dividing by $\sqrt{\det U_\mu^{\mathrm{new}}}$.  The action change is
	\begin{equation}\label{eq:deltaS}
		\Delta S = -\frac{\beta}{4}\,\Re\Tr\!\bigl[(U^{\mathrm{new}}-U^{\mathrm{old}})
		\cdot A^\dagger_\mu(x)\bigr],
	\end{equation}
	where $A_\mu(x)=\sum_{\nu\neq\mu}[A^{\mathrm{upper}}_{\mu\nu}(x)
	+A^{\mathrm{lower}}_{\mu\nu}(x)]$ is the sum of all six staples for link
	$(x,\mu)$, with
	\begin{align}
		A^{\mathrm{upper}}_{\mu\nu} &= U_\nu(x+\hat\mu)\,U_\mu^\dagger(x+\hat\nu)\,
		U_\nu^\dagger(x),\\
		A^{\mathrm{lower}}_{\mu\nu} &= U_\nu^\dagger(x+\hat\mu-\hat\nu)\,
		U_\mu^\dagger(x-\hat\nu)\,U_\nu(x-\hat\nu).
	\end{align}
	The proposal is accepted with probability $\min(1,e^{-\Delta S})$.  One full sweep
	proceeds through both parity classes and all four directions.  Thermalization used
	$N_{\mathrm{therm}}=500$ sweeps from a cold start (all links $=I_2$), followed by
	production with $N_{\mathrm{decorr}}=10$ decorrelation sweeps between measurements
	and $N_{\mathrm{cfg}}=50$ independent configurations per $\beta$.
	
	The step size $\varepsilon$ is auto-tuned before production to bring the
	acceptance rate to ${\sim}50\%$; the tuned values are $\varepsilon=0.50$ at
	$\beta=4.5$ (acceptance $0.485$) and $\varepsilon=0.35$ at $\beta=6.0$
	(acceptance $0.527$).  The plaquette observables are stable across configurations
	($\sigma_P\approx0.001$), indicating adequate equilibration.  The central result of
	this paper --- the channel equality $G_W^{(+)}=G_W^{(-)}$ --- is a per-configuration
	algebraic identity and is therefore independent of the acceptance rate or the degree
	of thermalization.

	
	\subsection{The SU(2) Link Embedding}
	
	The core methodological element is the embedding of $2\times2$ SU(2) link matrices
	into the 8-dimensional $\Cl(6)$ space.  The embedding is implemented as a
	vectorised NumPy operation using the Kronecker product identity:
	\begin{equation}
		\embed_{\SU(2)}(U) = \Id{2}\otimes U\otimes\Id{2},
	\end{equation}
	producing an $8\times8$ unitary matrix.  The analytical reduction
	\begin{equation}\label{eq:reduction}
		\Tr\!\bigl[P_r\,\embed(U_t)\,P_r\,\embed(U_0)^\dagger\bigr]
		= 4\cdot(U_t)_{rr}\cdot\overline{(U_0)_{rr}}
	\end{equation}
	is confirmed to machine precision for 50 random SU(2) pairs, where $r=0$ for
	$\Pew{+}$ and $r=1$ for $\Pew{-}$.
	
	\subsection{Numerical Environment}
	
	All computations are performed in Python~3.11 with NumPy~2.2, using IEEE~754
	double-precision arithmetic (machine epsilon
	$\varepsilon_{\mathrm{mach}}\approx2.2\times10^{-16}$).  Vectorised Kronecker
	products, \texttt{einsum}, and batched matrix multiplication via
	\texttt{einsum('nij,njk->nik')} are used throughout.
	
	% ══════════════════════════════════════════════════════════════════════════════
	\section{The Grade-Decomposed Link Correlator}
	
	\subsection{Definition}
	
	The grade-decomposed link correlator is defined by analogy with the
	grade-projected pion correlator of Paper~2~\cite{P2}, replacing the quark
	propagator (fermion bilinear) with the gauge link as the elementary correlating
	object.  For each gauge configuration and temporal separation~$t$, the projected
	correlator in the $r$-th EW grade is
	\begin{equation}
		\GW{r}(t) = \sum_{x_{\mathrm{sp}},\,\mu}
		\Re\Tr\!\bigl[\Pew{r}\,E_\mu(x,t)\,\Pew{r}\,E_\mu(x,0)^\dagger\bigr],
	\end{equation}
	where the sum runs over all spatial sites $x_{\mathrm{sp}}=(x_1,x_2,x_3)$ at
	fixed~$t$, over all four directions~$\mu$, and
	$E_\mu(x,t)=\Id{2}\otimes U_\mu(x,t)\otimes\Id{2}$ is the embedded link.  The
	total (unprojected) correlator is
	\begin{equation}
		G_W(t) = \sum_{x_{\mathrm{sp}},\,\mu}
		\Re\Tr\!\bigl[E_\mu(x,t)\,E_\mu(x,0)^\dagger\bigr].
	\end{equation}
	The grade fractions are
	\begin{equation}
		\fEW{r}(t) = \GW{r}(t)\,/\,G_W(t),\qquad r\in\{+,-\}.
	\end{equation}
	
	\subsection{Analytical Reduction}
	
	A key simplification follows from the structure of the embedding and the projectors.
	Using Eq.~(\ref{eq:reduction}), the correlators reduce to
	\begin{align}
		\GW{+}(t) &= 4\sum_{x_{\mathrm{sp}},\mu}
		\Re\!\bigl[(U_\mu(x,t))_{00}\,\overline{(U_\mu(x,0))_{00}}\bigr],\\
		\GW{-}(t) &= 4\sum_{x_{\mathrm{sp}},\mu}
		\Re\!\bigl[(U_\mu(x,t))_{11}\,\overline{(U_\mu(x,0))_{11}}\bigr].
	\end{align}
	Only the $(r,r)$ diagonal element of each $2\times2$ SU(2) matrix enters; the
	factor of~4 arises from the trace over the outer $\CC^2\otimes\CC^2$ factor.
	
	\subsection{The Channel Decomposition and Its Cross Term}
	
	The completeness relation $\Pew{+}+\Pew{-}=\Id{8}$ (check~B1) gives the exact
	per-configuration partition
	\begin{equation}\label{eq:partition}
		G_W(t) = \GW{+}(t)+\GW{-}(t)+\Delta(t),
	\end{equation}
	where the cross term collects the off-diagonal link contributions retained in the
	unprojected total but removed by the channel projectors,
	\begin{equation}\label{eq:crossdef}
		\Delta(t) = 8\sum_{x_{\mathrm{sp}},\mu}\Re\!\bigl[(U_\mu(x,t))_{01}\,
		\overline{(U_\mu(x,0))_{01}}\bigr].
	\end{equation}
	In the free field every link is $I_2$, so the off-diagonal element vanishes and
	$\Delta(t)=0$ identically; this is the special case underlying the free-field
	theorem below.  In the interacting theory $\Delta(t)$ is generically nonzero and
	$O(G_W)$: on our thermalized ensembles $|\Delta(t)|/G_W(t)$ reaches order unity.
	The partition~(\ref{eq:partition}) is an algebraic identity, verified per
	configuration (check~C6).
	
	\subsection{The Central Identity and the Free-Field Fractions}
	
	We first record the free-field value of the fractions, then state the exact
	identity that survives in the interacting theory.
	
	\begin{theorem}[Free-field fractions]\label{thm:freefield}
		In the free $\SU(2)$ gauge theory (all links $U_\mu=I_2$), the channel fractions
		$\fEW{\pm}(t)=\GW{\pm}(t)/G_W(t)$ equal $1/2$ for both $r\in\{+,-\}$ and all
		timeslices.
	\end{theorem}
	
	\begin{proof}
		With all links $=I_2$, $\embed(I_2)=\Id{8}$, so
		\[
		\GW{r}(t)
		= \sum_{x_{\mathrm{sp}},\mu}\Tr\!\bigl[(\Pew{r})^2\bigr]
		= \sum_{x_{\mathrm{sp}},\mu}\Tr[\Pew{r}]
		= 4\cdot N^3\cdot 4 = 16N^3,
		\]
		\[
		G_W(t) = \sum_{x_{\mathrm{sp}},\mu}\Tr[\Id{8}] = 8\cdot N^3\cdot 4 = 32N^3,
		\]
		giving $\fEW{r}(t)=16N^3/32N^3=1/2$ for both $r$.  This rests on $\Delta=0$ in
		the free field; it does \emph{not} extend to the interacting theory, where
		$\Delta\neq0$ (Theorem~\ref{thm:channel} and \S\ref{sec:gauge}).\qed
	\end{proof}
	
	The central result of the paper is a different statement, which holds at all
	couplings.
	
	\begin{theorem}[Per-configuration channel equality]\label{thm:channel}
		For every $\SU(2)$ gauge configuration and every timeslice,
		$\GW{+}(t)=\GW{-}(t)$.
	\end{theorem}
	
	\begin{proof}
		Any $U\in\SU(2)$ has diagonal elements satisfying $U_{11}=\overline{U_{00}}$;
		writing $U_{00}=a$, we have $U_{11}=\bar a$, with $|a|^2+|U_{01}|^2=1$.  By the
		reduction~(\ref{eq:reduction}), $\GW{+}\propto\sum\Re[a_t\bar a_0]$ and
		$\GW{-}\propto\sum\Re[\bar a_t a_0]$.  Since $\bar a_t a_0=\overline{a_t\bar a_0}$
		and $\Re z=\Re\bar z$, the two sums are identical term by term.  Hence
		$\GW{+}(t)=\GW{-}(t)$ for every configuration.\qed
	\end{proof}
	
	\begin{remark}[Origin: pseudoreality]
		Theorem~\ref{thm:channel} is the lattice-correlator expression of the
		pseudoreality of SU(2): the reality structure $\varepsilon\,\overline{U}\,
		\varepsilon^{-1}=U$ with $\varepsilon=i\sigma_y$ exchanges the two $T_3$
		eigen-channels, forcing their diagonal correlators to coincide.  The result is
		therefore structural --- a property of the group, made manifest by the $T_3$
		grading --- rather than a dynamical feature of any particular ensemble.
	\end{remark}
	
	\begin{remark}[Why this does not give $\fEW{\pm}=1/2$]
		A natural but incorrect extrapolation from Theorem~\ref{thm:freefield} would be
		that $\fEW{\pm}=1/2$ at all couplings.  Theorem~\ref{thm:channel} gives only
		$\GW{+}=\GW{-}$, hence $\fEW{+}=\fEW{-}$ --- the two fractions are equal to each
		other, not to $1/2$.  Equality with $1/2$ additionally requires
		$\GW{+}+\GW{-}=G_W$, i.e.\ $\Delta=0$, which holds only in the free field.  Since
		$\Delta\neq0$ in the interacting theory~(\ref{eq:crossdef}), the fractions are not
		pinned to $1/2$; this is confirmed in \S\ref{sec:results}.
	\end{remark}
	
	\subsection{Gauge Dependence of the Channel Correlators}
	\label{sec:gauge}
	
	The link correlator of this paper is built from two links at the same spatial point
	and direction but different times, $E_\mu(x,t)$ and $E_\mu(x,0)$, with no connecting
	Wilson line. Such a product is not gauge invariant: under a gauge transformation
	$U_\mu(x)\to g(x)\,U_\mu(x)\,g(x+\hat\mu)^\dagger$, the open product
	$U_\mu(x,t)U_\mu(x,0)^\dagger$ transforms with uncancelled factors at the endpoints.
	Consequently $G_W(t)$, $\GW{\pm}(t)$, and $\Delta(t)$ are all gauge-variant for
	$t\neq0$.
	
	By Elitzur's theorem~\cite{Elitzur1975}, the ensemble average of a gauge-non-invariant
	local operator vanishes in the absence of gauge fixing. The expectation values
	$\langle G_W(t)\rangle$ and $\langle\GW{\pm}(t)\rangle$ are therefore zero (up to
	finite-ensemble noise) for $t\neq0$, and the \emph{averaged} fractions
	$\langle\GW{\pm}\rangle/\langle G_W\rangle$ are ratios of quantities consistent with
	zero --- numerically ill-defined, with errors that do not shrink in any controlled
	way. We do not present ensemble-averaged channel correlators or fractions as physical
	observables.
	
	This is the reason the result of this paper is stated as a \emph{per-configuration}
	identity. Theorem~\ref{thm:channel} holds for each gauge configuration separately ---
	it is a statement about the integrand, not the integral --- and is therefore
	unaffected by Elitzur's theorem, gauge fixing, or thermalization. The equality
	$\GW{+}(t)=\GW{-}(t)$ is well-posed configuration by configuration regardless of
	whether the individual correlators have a nonzero average. A gauge-invariant
	formulation (for example, connecting the two links by a Wilson line, or fixing to a
	smooth gauge) would be required before any ensemble-averaged channel quantity could
	be interpreted physically; we leave this to future work.
	
	% ══════════════════════════════════════════════════════════════════════════════
	\section{Results}
	
\subsection{Free-Field Verification}
	\label{sec:results}
	
	Setting all links to $I_2$ gives $\GW{\pm}(t)=\Tr[\Pew{\pm}]\cdot N^3\cdot4=16N^3$
	and $G_W(t)=32N^3$, so $\fEW{\pm}(t)=1/2$ exactly (Theorem~\ref{thm:freefield}).
	The simulation reproduces this with maximum deviation $0.00\times10^0$ for all $t$
	and the free-field gate checks (C1--C6) pass identically (Appendix~A).
	
	\subsection{Plaquette Equilibration}
	
	The ensembles equilibrate to stable plaquette values consistent with the standard
	SU(2) Wilson benchmarks at the corresponding $\beta_{\mathrm{std}}=\beta/2$
	(Table~\ref{tab:plaquette}).
	
	\begin{table}[htbp]
		\centering
		\caption{SU(2) plaquette equilibration. $\beta$ is in the action normalization of
			\S\ref{sec:lattice}; $\beta_{\mathrm{std}}=\beta/2$ is the standard Wilson
			coupling.}
		\label{tab:plaquette}
		\begin{tabular}{cccccc}
			\toprule
			$\beta$ & $\beta_{\mathrm{std}}$ & $\langle P\rangle$ & $\sigma_P$ &
			$\varepsilon$ & Mean acc.\\
			\midrule
			4.5 & 2.25 & 0.576243 & 0.000744 & 0.50 & 0.485\\
			6.0 & 3.00 & 0.718868 & 0.000535 & 0.35 & 0.527\\
			\bottomrule
		\end{tabular}
	\end{table}
	
	\noindent These values match standard SU(2) results: $\langle P\rangle\approx0.576$
	at $\beta_{\mathrm{std}}=2.25$ (just below deconfinement) and $\approx0.719$ at
	$\beta_{\mathrm{std}}=3.0$ (deconfined), confirming correct sampling.
	
	\subsection{The Partition Identity and the Size of the Cross Term}
	
	The per-configuration partition $G_W=\GW{+}+\GW{-}+\Delta$~(\ref{eq:partition})
	holds to machine precision on every configuration (Table~\ref{tab:completeness}).
	In contrast to the free field, the cross term $\Delta$ is \emph{not} negligible on
	thermalized ensembles: $|\Delta(t)|/G_W(t)$ reaches order unity, which is precisely
	why the channel fractions are not pinned to $1/2$ (\S\ref{sec:nofraction}).
	
	\begin{table}[htbp]
		\centering
		\caption{Per-configuration partition identity and cross-term size.}
		\label{tab:completeness}
		\begin{tabular}{ccc}
			\toprule
			$\beta$ &
			$\max_t\bigl|G_W-\GW{+}-\GW{-}-\Delta\bigr|/|G_W|$ &
			$\max_t|\Delta(t)|/|G_W(t)|$\\
			\midrule
			4.5 & $<10^{-15}$ (PASS) & 1.27\\
			6.0 & $<10^{-15}$ (PASS) & 2.21\\
			\bottomrule
		\end{tabular}
	\end{table}
	
	\subsection{The Channel Equality --- Primary Result}
	
	The central result, Theorem~\ref{thm:channel}, is the per-configuration equality
	$\GW{+}(t)=\GW{-}(t)$. On the thermalized ensembles this holds to machine precision
	on every configuration and every timeslice (Table~\ref{tab:channel}); the maximum
	relative difference is $0.00\times10^0$ (an exact equality of floating-point sums,
	since the two channels are identical term by term). This confirms that the identity
	is not an artifact of special configurations: it holds across both sampled phases.
	
	\begin{table}[htbp]
		\centering
		\caption{Per-configuration channel equality, verified over $N_{\mathrm{cfg}}=50$
			configurations at each coupling.}
		\label{tab:channel}
		\begin{tabular}{ccc}
			\toprule
			$\beta$ &
			$\max_{t,\,\mathrm{cfg}}\bigl|\GW{+}(t)-\GW{-}(t)\bigr|/|\GW{+}(t)|$ & Status\\
			\midrule
			4.5 & 0.00e+00 & PASS\\
			6.0 & 0.00e+00 & PASS\\
			\bottomrule
		\end{tabular}
	\end{table}
	
	\noindent We deliberately do \emph{not} tabulate ensemble-averaged fractions
	$\langle\GW{\pm}\rangle/\langle G_W\rangle$: as established in \S\ref{sec:gauge},
	these are gauge-variant ratios whose averages are consistent with zero and carry no
	physical meaning.
	
	\subsection{Comparison with the Color Sector}
	
	Both the color sector (Paper~2) and the present electroweak sector admit a grade
	decomposition of a lattice correlator into projector channels. The structural
	distinction is in what the decomposition makes exact. The SU(3) color projectors of
	Paper~2 have unequal dimensions ($3,3,1,1$) and there is no group-reality relation
	forcing distinct channels to coincide; the color-grade fractions take their
	free-field values only approximately on interacting ensembles.
	
	The SU(2) $T_3$ decomposition is different in kind. The two channels have equal
	dimension ($4+4$), and SU(2) pseudoreality forces them to be \emph{exactly equal
	per configuration}, $\GW{+}=\GW{-}$ (Theorem~\ref{thm:channel}), at all couplings.
	This per-configuration channel equality has no SU(3) analogue --- it is a
	consequence of the pseudoreality of the $2$ of SU(2), a property the $3$ of SU(3)
	does not share. We stress that this is a structural identity of the decomposition,
	not a statement about ensemble-averaged observables (\S\ref{sec:gauge}), and in
	particular not a ``kinematic theorem'' for the channel fractions, which are pinned
	to $1/2$ only in the free field.
	
	% ══════════════════════════════════════════════════════════════════════════════
\section{Discussion}
	
	\subsection{What the Identity Is, and What It Is Not}
	\label{sec:nofraction}
	
	The result of this paper is the per-configuration channel equality
	$\GW{+}(t)=\GW{-}(t)$ (Theorem~\ref{thm:channel}), an exact consequence of SU(2)
	pseudoreality. It is worth being precise about its status, because a weaker and a
	stronger reading are easy to confuse.
	
	The identity is \emph{structural}: it states that the $T_3$ grade decomposition of
	the SU(2) link correlator splits it into two channels that the group's reality
	structure forces to coincide. It holds for every configuration, at every coupling,
	in either phase, with or without gauge fixing. What it does \emph{not} provide is
	new dynamical information. The equality of the two channels is equivalent to the
	textbook fact that the $2$ of SU(2) is pseudoreal; the Clifford-algebra framework
	makes the split natural and the equality manifest, but does not extract a quantity
	inaccessible to the standard formulation.
	
	In particular, the identity must not be read as the statement that the channel
	fractions are pinned to $1/2$. That would require the cross term $\Delta$ to vanish,
	which happens only in the free field (Theorem~\ref{thm:freefield}). On interacting
	ensembles $\Delta$ is $O(G_W)$ (Table~\ref{tab:completeness}), so the fractions are
	not $1/2$; and in any case the fractions, being built from gauge-variant correlators,
	are not physical observables (\S\ref{sec:gauge}). The surviving, well-posed content
	is the channel equality alone.
	
	\subsection{Relation to a Natural Prior Expectation}
	
	It is instructive to record the expectation this calculation corrects. The free-field
	result $\fEW{\pm}=1/2$ (Theorem~\ref{thm:freefield}), together with the exact
	per-configuration equality $\GW{+}=\GW{-}$, invites the conclusion that
	$\fEW{\pm}=1/2$ holds at all couplings. This is false. Equality of the two channels
	does not pin either fraction to $1/2$; only the additional condition $\Delta=0$,
	special to the free field, does. The interacting cross term is not a small
	correction --- it is order unity --- so the free-field fraction does not survive the
	introduction of dynamics. The robust statement is the channel equality, not the
	fraction value.
	
	\subsection{The \texorpdfstring{$\beta$}{beta} Regime}
	
	In the action normalization of \S\ref{sec:lattice}, the SU(2) deconfinement
	transition ($\beta_{\mathrm{std}}\approx2.30$) sits at $\beta\approx4.6$. The two
	couplings studied, $\beta=4.5$ and $\beta=6.0$, lie just below and just above the
	transition respectively, and the plaquette values (Table~\ref{tab:plaquette})
	confirm the corresponding confined and deconfined behavior. The channel equality
	holds identically in both, as it must: being a per-configuration identity, it is
	insensitive to the phase. No claim about confinement or symmetry breaking is made
	or implied.
	
	\subsection{What Cannot Be Concluded}
	
	This computation does not address, and should not be interpreted as evidence for or
	against:
	
	\begin{enumerate}[leftmargin=*,label=\arabic*.]
		\item \textbf{The Higgs mechanism.}  Pure gauge theory without scalars; the
		Derivability Limit (\S2.4) applies.
		\item \textbf{W and Z boson masses.}  Mass extraction requires a gauge-invariant
		correlator and a state interpretation not available here.
		\item \textbf{Chiral structure.}  Nielsen--Ninomiya~\cite{NN1,NN2} prohibits exact
		chiral symmetry on the lattice.
		\item \textbf{L/R handedness.}  SU(2) pure gauge theory is vector-like; L/R
		assignment requires coupling to left-handed matter.
		\item \textbf{Any ensemble-averaged channel quantity.}  The averaged correlators
		and fractions are gauge-variant and vanish by Elitzur's theorem (\S\ref{sec:gauge}).
	\end{enumerate}
	
	\subsection{Outlook}
	
	The most natural way to give the channel decomposition dynamical content is to make
	it gauge invariant. Two routes are available: connecting the two links by a Wilson
	line so that $\GW{\pm}$ become gauge-invariant correlators, or fixing to a smooth
	gauge (Landau or Coulomb) and studying the residual-gauge-invariant content. In
	either case one could ask whether any structure beyond the pseudoreality equality
	survives. A second direction is to embed the decomposition in the matter sector,
	where the channels would act on fermion bilinears rather than bare links; there the
	analogue of $\Delta$ may carry information that the pure-gauge, gauge-variant version
	cannot. Both are left to future work.
	
	% ══════════════════════════════════════════════════════════════════════════════
	\section{Conclusions}
	
	This paper has studied the $T_3$ grade decomposition of the SU(2) link correlator in
	the $\Cl(2)\cong\HH$ framework of this series, establishing the following.
	
	\textbf{(i) Framework definition and algebraic verification.}  The grade-decomposed
	link correlator is defined by embedding SU(2) links into the 8-dimensional
	representation space via $U_\mu^{(8)}=\Id{2}\otimes U_\mu\otimes\Id{2}$.  The
	generator algebra, the $T_3$ spectrum, and the channel projectors (idempotence,
	completeness, orthogonality, hermiticity, equal trace, embedding unitarity, and the
	analytical reduction of the projected correlator) all pass at machine precision
	(Table~\ref{tab:algebra}).
	
	\textbf{(ii) The central result: per-configuration channel equality.}  For every
	SU(2) gauge configuration the two $T_3$ channel correlators are exactly equal,
	$\GW{+}(t)=\GW{-}(t)$ (Theorem~\ref{thm:channel}), verified to $0.00\times10^0$ on
	thermalized ensembles at $\beta\in\{4.5,6.0\}$ in both the confined and deconfined
	phases.  The identity follows from the SU(2) diagonal relation
	$U_{11}=\overline{U_{00}}$ --- the pseudoreality of the group --- and holds per
	configuration, independent of coupling, phase, gauge fixing, or thermalization.
	
	\textbf{(iii) Honest scope of the result.}  The equality is a structural feature of
	the decomposition, not new dynamical information: it is the lattice expression of
	SU(2) pseudoreality.  The channel fractions equal $1/2$ only in the free field
	(Theorem~\ref{thm:freefield}); in the interacting theory the cross term is order
	unity and the fractions are not pinned to $1/2$.  Moreover the channel correlators
	are gauge-variant, so their ensemble averages are not physical observables
	(Elitzur's theorem, \S\ref{sec:gauge}).  The well-posed content is the
	per-configuration channel equality alone.
	
	\textbf{What is not established.}  This paper does not test the Higgs mechanism,
	W- and Z-boson masses, chiral structure, or SU(2)$_L$ vs.\ SU(2)$_R$ handedness, and
	makes no claim about any ensemble-averaged channel quantity.  A gauge-invariant
	formulation of the decomposition is the natural next step.
	
	% ══════════════════════════════════════════════════════════════════════════════
	\appendix
	
	\section{SU(2) Algebraic Verification Table}
	
	All checks are computed at the start of the simulation from the exact matrix
	definitions of this paper.  Checks marked \emph{analytic} are verified by
	algebraic argument; all others are code-verified.  Tolerances $<10^{-12}$ are
	considered machine precision.
	
	\begin{table}[htbp]
		\centering
		\caption{Algebraic verification checks}
		\label{tab:algebra}
		\begin{tabular}{clccc}
			\toprule
			Check & Description & Value & Method & Status\\
			\midrule
			A1 & $[T_i,T_j]=i\varepsilon_{ijk}T_k$ (all three) & 0.000e+00 & code & PASS\\
			A2 & $T^2=(3/4)\Id{8}$                              & 0.000e+00 & code & PASS\\
			A3 & $T_3$ eigenvalues $\in\{\pm1/2\}\times4$       & 0.000e+00 & code & PASS\\
			B1 & $\Pew{+}+\Pew{-}=\Id{8}$                       & 0.000e+00 & code & PASS\\
			B2 & $\Pew{+}\Pew{-}=0$                              & 0.000e+00 & code & PASS\\
			B3 & $(\Pew{\pm})^2=\Pew{\pm}$                       & 0.000e+00 & code & PASS\\
			B4 & $\Pew{\pm}$ Hermitian                           & ---       & analytic & PASS\\
			B5 & $\dim(\Pew{\pm})=4$ each                        & 0.000e+00 & code & PASS\\
			B6 & $\embed(I_2)=\Id{8}$                            & 0.000e+00 & code & PASS\\
			B7 & $\|\embed(U)\embed(U)^\dagger-\Id{8}\|_F$
			(20 random $U$)                                  & 1.26e-15  & code & PASS\\
			B8 & $[P_{\mathrm{color}},\Pew{}]=0$                 & ---       & analytic & PASS\\
			V9  & $\GW{+}-\GW{-}=0$ per config (max)             & 6.30e-16  & code & PASS\\
			V10 & Eq.~(\ref{eq:reduction}) analytical reduction & 4.44e-16  & code & PASS\\
			\bottomrule
		\end{tabular}
	\end{table}
	
	\section{Simulation Parameters}
	
	\begin{table}[htbp]
		\centering
		\caption{Simulation parameters}
		\label{tab:simparams}
		\begin{tabular}{lll}
			\toprule
			Parameter & $\beta=4.5$ & $\beta=6.0$\\
			\midrule
			Gauge group                                    & \multicolumn{2}{l}{SU(2)}\\
			Spacetime dimension                            & \multicolumn{2}{l}{4}\\
			Lattice size $N$                               & \multicolumn{2}{l}{6}\\
			Volume $V=N^4$                                 & \multicolumn{2}{l}{1296}\\
			$\beta_{\mathrm{std}}=\beta/2$                 & 2.25 & 3.00\\
			Action prefactor $\Delta S$                    & \multicolumn{2}{l}{$\beta/4$ per link}\\
			$\varepsilon$ (auto-tuned Metropolis step)     & 0.50 & 0.35\\
			Sweep algorithm                                & \multicolumn{2}{l}{Even-odd checkerboard}\\
			Start configuration                            & \multicolumn{2}{l}{Cold (all links $=I_2$)}\\
			Thermalization sweeps $N_{\mathrm{therm}}$     & \multicolumn{2}{l}{500}\\
			Decorrelation sweeps $N_{\mathrm{decorr}}$     & \multicolumn{2}{l}{10}\\
			Configurations per $\beta$                     & \multicolumn{2}{l}{50}\\
			Acceptance rate                                & 48.5\% & 52.7\%\\
			$\langle P\rangle$                             & $0.5762\pm0.0007$ & $0.7189\pm0.0005$\\
			Statistical errors                             & \multicolumn{2}{l}{Jackknife ($N=50$)}\\
			Language / library                             & \multicolumn{2}{l}{Python~3.11 / NumPy~2.2.5}\\
			Arithmetic                                     & \multicolumn{2}{l}{IEEE~754 double precision}\\
			$\varepsilon_{\mathrm{mach}}$                  & \multicolumn{2}{l}{$\approx2.2\times10^{-16}$}\\
			\bottomrule
		\end{tabular}
	\end{table}
	
	% ══════════════════════════════════════════════════════════════════════════════
	\begin{thebibliography}{99}
		
		\bibitem{P1}
		[Author~et~al.], ``The Division Algebra Framework for the Standard Model:
		$\RR\otimes\CC\otimes\HH\otimes\OO$ and $\Cl(6)$,''
		\textit{Adv.\ Appl.\ Clifford Algebras}, [year].
		(Paper~1 of this series)
		
		\bibitem{P2}
		[Author~et~al.], ``Grade Decomposition of the Lattice QCD Pion Correlator
		in the $\Cl(6)$ Framework,''
		\textit{Adv.\ Appl.\ Clifford Algebras}, [year].
		(Paper~2 of this series)
		
		\bibitem{P3}
		[Author~et~al.], ``SU(2)$_L$, the Higgs Doublet, and Electroweak Symmetry
		Breaking in $\RR\otimes\CC\otimes\HH\otimes\OO$,''
		\textit{Adv.\ Appl.\ Clifford Algebras}, [year].
		(Paper~3 of this series)
		
		\bibitem{PS1974}
		J.\,C.~Pati and A.~Salam,
		``Lepton Number as the Fourth Color,''
		\textit{Phys.\ Rev.\ D} \textbf{10} (1974) 275--289.
		
		\bibitem{S1981}
		R.~Slansky,
		``Group Theory for Unified Model Building,''
		\textit{Phys.\ Rep.} \textbf{79} (1981) 1--128.
		
		\bibitem{NN1}
		H.\,B.~Nielsen and M.~Ninomiya,
		``Absence of Neutrinos on a Lattice: (I).\ Proof by Homotopy Theory,''
		\textit{Nucl.\ Phys.\ B} \textbf{185} (1981) 20--40.
		
		\bibitem{NN2}
		H.\,B.~Nielsen and M.~Ninomiya,
		``Absence of Neutrinos on a Lattice: (II).\ Intuitive Topological Proof,''
		\textit{Nucl.\ Phys.\ B} \textbf{193} (1981) 173--194.
		
		\bibitem{L1998}
		M.~L\"uscher,
		``Advanced Lattice QCD,''
		\textit{Les Houches Lectures}, 1998, hep-lat/9802029.
		
		\bibitem{K1979}
		J.~Kogut, R.\,B.~Pearson, and J.~Shigemitsu,
		``The Quark-Confining String and the World Sheet,''
		\textit{Phys.\ Rev.\ Lett.} \textbf{43} (1979) 484--486.
		
		\bibitem{BC1980}
		G.~Bhanot and M.~Creutz,
		``Variant Actions and Phase Structure in Lattice Gauge Theory,''
		\textit{Phys.\ Rev.\ D} \textbf{24} (1981) 3212--3217.
		
		\bibitem{Elitzur1975}
		S.~Elitzur,
		``Impossibility of spontaneously breaking local symmetries,''
		\textit{Phys.\ Rev.\ D} \textbf{12} (1975) 3978--3982.
		
	\end{thebibliography}
	
\end{document}